\section{Justificación}
En este contexto, el presente proyecto de trabajo de grado busca aportar a la comprensión de la inestabilidad de Kelvin–Helmholtz mediante una caracterización numérica detallada en el marco de la Magnetohidrodinámica Relativista Resistiva (RRMHD). Este enfoque permite analizar con mayor rigor cómo parámetros físicos fundamentales, como la resistividad, el contraste de densidad y la orientación del campo magnético, influyen en el desarrollo y la evolución tanto lineal como no lineal de la inestabilidad. Al abordar explícitamente estos efectos no ideales, el estudio pretende capturar dinámicas que no pueden describirse adecuadamente mediante modelos hidrodinámicos o magnetohidrodinámicos ideales, contribuyendo así a una comprensión más completa de los procesos que gobiernan la estabilidad de plasmas en entornos astrofísicos extremos.

Sobre esta base, y una vez definidos los parámetros que condicionan la evolución de la inestabilidad, el trabajo se orienta a explorar a  través de simulaciones numéricas controladas, los mecanismos que gobiernan el comportamiento no lineal de esta inestabilidad en condiciones astrofísicas extremas. Los resultados obtenidos podrían ser de gran utilidad para fortalecer y validar modelos más generales dentro del campo de la astrofísica relativista, especialmente en el estudio de chorros relativistas, magnetósferas estelares y discos de acreción.